\scanpage{42}
des groupes \inserted{profinis} extérieurs avec opérations de $\Gamma$ dessus. Et on pressent que le th. 3 du \S\ précédent (appliqué notamment à $\mathbb P_K^1$ convenablement troué...) pourrait donner la clef d'une telle construction.

Bien \deleted{\unsure{en}} sûr, des homomorphismes extérieurs quelconques $\pi_{\overline{K'}}\longrightarrow\pi_{\overline K}$ n'auront pas de sens géométrique --- l'idée est que les opérations du groupe $\Gamma=\Gamma_{\overline{\mathbb Q}/k}$ dessus sont si draconiennes, qu'il n'est possible de trouver un hom. extérieur qui y commute que par voie géométrique --- par des plongements de corps. \unsure{Pour|Or} il est essentiel ici que le corps de base ne soit pas quelconque, mais un corps tel que $\mathbb Q$ (ou, ce qui revient manifestement au $\unsure{\overline m}$, une extension de type fini de $\mathbb Q$). Encore faut-il se borner aux hom. $\pi_{K'}-\pi_K$ dont on décrète d'avance que l'image soit ouverte --- sinon
\margin{41}
